\documentclass[11pt,a4paper]{article}

% Packages essentiels
\usepackage[utf8]{inputenc}
\usepackage[french]{babel}
\usepackage[T1]{fontenc}
\usepackage{amsmath}
\usepackage{amssymb}
\usepackage{graphicx}
\usepackage{geometry}
\usepackage{hyperref}
\usepackage{enumitem}
\usepackage{booktabs}
\usepackage{xcolor}
\usepackage{fancyhdr}
\usepackage{titlesec}
\usepackage{listings}
\usepackage{tcolorbox}
\usepackage{array}
\usepackage{multirow}

% Configuration de la page
\geometry{margin=2.5cm}

% Couleurs personnalisées
\definecolor{darkblue}{RGB}{0,51,102}
\definecolor{lightblue}{RGB}{59,130,246}
\definecolor{codegreen}{RGB}{0,128,0}
\definecolor{codegray}{RGB}{128,128,128}
\definecolor{codepurple}{RGB}{128,0,128}
\definecolor{backcolour}{RGB}{245,245,245}

% Configuration des en-têtes et pieds de page
\pagestyle{fancy}
\fancyhf{}
\fancyhead[L]{\small Application de Distillation Multicomposants}
\fancyhead[R]{\small Rapport Complet v2.2}
\fancyfoot[C]{\thepage}

% Formatage des titres
\titleformat{\section}{\Large\bfseries\color{darkblue}}{\thesection}{1em}{}
\titleformat{\subsection}{\large\bfseries\color{lightblue}}{\thesubsection}{1em}{}

% Configuration des liens hypertexte
\hypersetup{
    colorlinks=true,
    linkcolor=darkblue,
    urlcolor=lightblue,
    citecolor=darkblue,
    pdfborder={0 0 0}
}

% Configuration pour le code source
\lstdefinestyle{pythonstyle}{
    backgroundcolor=\color{backcolour},
    commentstyle=\color{codegreen},
    keywordstyle=\color{blue},
    numberstyle=\tiny\color{codegray},
    stringstyle=\color{codepurple},
    basicstyle=\ttfamily\footnotesize,
    breakatwhitespace=false,
    breaklines=true,
    captionpos=b,
    keepspaces=true,
    numbers=left,
    numbersep=5pt,
    showspaces=false,
    showstringspaces=false,
    showtabs=false,
    tabsize=2,
    language=Python
}

\lstset{style=pythonstyle}

% Commandes personnalisées
\newcommand{\code}[1]{\texttt{#1}}
\newcommand{\file}[1]{\texttt{\textbf{#1}}}

% Début du document
\begin{document}

% ============================================================================
% PAGE DE TITRE
% ============================================================================
\begin{titlepage}
    \centering
    \vspace*{1.5cm}

    {\Huge\bfseries Application de Distillation\\Vspace{0.5cm}Multicomposants\par}
    \vspace{1cm}

    {\Large\bfseries Rapport Complet de l'Application\par}
    \vspace{0.5cm}
    {\large Documentation Technique et Guide d'Utilisation\par}

    \vspace{2cm}

    \begin{tcolorbox}[colback=lightblue!10,colframe=darkblue,title=Informations du Projet]
        \textbf{Module:} Modélisation et Simulation des Procédés\\[0.3cm]
        \textbf{Professeur:} BAKHER Zine Elabidine\\[0.3cm]
        \textbf{Filière:} Procédés Industriels et Chimiques (PIC)\\[0.3cm]
        \textbf{Université:} Université Hassan 1er\\[0.3cm]
        \textbf{Année Académique:} 2024-2025
    \end{tcolorbox}

    \vspace{2cm}

    \begin{tcolorbox}[colback=green!10,colframe=green!50!black,title=Version de l'Application]
        \textbf{Version:} 2.2 - Numérotation Séquentielle et Téléchargement PDF\\[0.3cm]
        \textbf{Date:} 27 Novembre 2025\\[0.3cm]
        \textbf{Statut:} \textcolor{green!70!black}{\textbf{PRODUCTION READY}}
    \end{tcolorbox}

    \vfill

    {\large \textit{La distillation - Procédés de Séparation}\par}
\end{titlepage}

% Table des matières
\tableofcontents
\newpage

% ============================================================================
\section{Vue d'Ensemble de l'Application}
% ============================================================================

\subsection{Introduction}

Cette application web interactive développée avec \textbf{Streamlit} constitue un outil pédagogique complet pour l'étude de la distillation multicomposants. Elle implémente l'intégralité du programme du cours "Modélisation et Simulation des Procédés" de la filière PIC à l'Université Hassan 1er.

\subsection{Objectifs Pédagogiques}

L'application vise à permettre aux étudiants de:

\begin{itemize}[leftmargin=*]
    \item \textbf{Comprendre} les principes fondamentaux de la distillation multicomposants
    \item \textbf{Appliquer} les méthodes simplifiées de dimensionnement (Fenske, Underwood, Gilliland, Kirkbride)
    \item \textbf{Utiliser} la méthode rigoureuse MESH avec l'algorithme Wang-Henke
    \item \textbf{Modéliser} les mélanges non-idéaux avec Wilson, NRTL et UNIQUAC
    \item \textbf{Optimiser} économiquement une colonne par minimisation du TAC
    \item \textbf{Visualiser} les profils de composition, température et débits
    \item \textbf{Analyser} les bilans matière et énergétiques
\end{itemize}

\subsection{Caractéristiques Principales}

\begin{tcolorbox}[colback=blue!5,colframe=blue!50!black,title=Fonctionnalités Clés]
\begin{itemize}[leftmargin=*]
    \item \textbf{Interface Intuitive:} Design moderne et professionnel avec navigation simplifiée
    \item \textbf{Bibliothèque de 13 Composés:} Hydrocarbures aromatiques, alcools, alcanes
    \item \textbf{Trois Modes de Calcul:} Méthodes Simplifiées, MESH Rigoureux, Comparaison
    \item \textbf{Visualisations Interactives:} Graphiques Plotly avec zoom, pan, export
    \item \textbf{Modèles Thermodynamiques:} Idéal, Wilson, NRTL, UNIQUAC
    \item \textbf{Optimisation Économique:} Calcul du TAC, études paramétriques
    \item \textbf{Documentation Intégrée:} 20 équations numérotées, théorie complète
    \item \textbf{Export PDF:} Téléchargement de la documentation LaTeX
\end{itemize}
\end{tcolorbox}

% ============================================================================
\section{Architecture de l'Application}
% ============================================================================

\subsection{Structure Modulaire}

L'application est organisée en modules Python indépendants pour faciliter la maintenance et l'évolution:

\begin{table}[h]
\centering
\small
\begin{tabular}{p{5cm}p{9cm}}
\toprule
\textbf{Module} & \textbf{Fonctionnalité} \\
\midrule
\file{streamlit\_app\_enhanced.py} & Interface utilisateur Streamlit (1500+ lignes) \\
\file{distillation\_multicomposants.py} & Moteur de calcul - méthodes simplifiées (800+ lignes) \\
\file{mesh\_solver.py} & Solveur MESH rigoureux - algorithme Wang-Henke (500+ lignes) \\
\file{activity\_models.py} & Modèles thermodynamiques non-idéaux (600+ lignes) \\
\file{economic\_optimization.py} & Optimisation économique TAC (500+ lignes) \\
\bottomrule
\end{tabular}
\caption{Modules Python de l'application (3900+ lignes de code au total)}
\end{table}

\subsection{Flux de Données}

Le diagramme de flux de calcul se déroule comme suit:

\begin{enumerate}[leftmargin=*]
    \item \textbf{Entrée Utilisateur} (Streamlit UI):
    \begin{itemize}
        \item Sélection des composés (2 à 6)
        \item Compositions de l'alimentation $z_i$
        \item Paramètres opératoires ($F$, $P$, $R$, $\eta$)
        \item Spécifications de séparation (récupérations)
    \end{itemize}

    \item \textbf{Méthodes Simplifiées} (\file{distillation\_multicomposants.py}):
    \begin{itemize}
        \item Fenske: $N_{min}$ (Éq. 1)
        \item Underwood: $R_{min}$, $\theta$ (Éq. 2-3)
        \item Gilliland: $N$ théorique (Éq. 4-6)
        \item Efficacité: $N_{real}$ (Éq. 7)
        \item Kirkbride: Position alimentation (Éq. 8)
    \end{itemize}

    \item \textbf{MESH Rigoureux} (optionnel, \file{mesh\_solver.py}):
    \begin{itemize}
        \item Initialisation avec résultats simplifiés
        \item Résolution itérative Wang-Henke (Éq. 9-13)
        \item Modèles d'activité si mélange non-idéal (Éq. 14-16)
        \item Convergence sur températures ($\epsilon = 10^{-6}$)
    \end{itemize}

    \item \textbf{Optimisation Économique} (\file{economic\_optimization.py}):
    \begin{itemize}
        \item Calcul des coûts d'investissement (Éq. 18)
        \item Calcul des coûts d'exploitation (Éq. 19)
        \item Total Annualized Cost - TAC (Éq. 17)
        \item Coût de maintenance (Éq. 20)
    \end{itemize}

    \item \textbf{Visualisation} (Plotly):
    \begin{itemize}
        \item Génération de graphiques interactifs
        \item Affichage des résultats dans l'interface
        \item Export des données et figures
    \end{itemize}
\end{enumerate}

\subsection{Technologies Utilisées}

\begin{table}[h]
\centering
\begin{tabular}{llp{6cm}}
\toprule
\textbf{Technologie} & \textbf{Version} & \textbf{Utilisation} \\
\midrule
Python & 3.8+ & Langage de programmation principal \\
Streamlit & Latest & Framework d'interface utilisateur web \\
NumPy & Latest & Calculs numériques et algèbre linéaire \\
SciPy & Latest & Optimisation et résolution d'équations \\
Plotly & Latest & Visualisations interactives \\
Pandas & Latest & Manipulation de données tabulaires \\
Thermo & Latest & Propriétés thermodynamiques \\
\bottomrule
\end{tabular}
\caption{Stack technologique de l'application}
\end{table}

% ============================================================================
\section{Fondements Théoriques Implémentés}
% ============================================================================

\subsection{Bilans Matière Globaux}

Les équations fondamentales implémentées dans tous les modes de calcul:

\begin{equation}
F = D + B
\end{equation}

\begin{equation}
F \cdot z_i = D \cdot x_{D,i} + B \cdot x_{B,i} \quad \forall i
\end{equation}

où:
\begin{itemize}[leftmargin=*]
    \item $F$ = Débit d'alimentation (kmol/h)
    \item $D$ = Débit de distillat (kmol/h)
    \item $B$ = Débit de résidu (kmol/h)
    \item $z_i$ = Fraction molaire du composé $i$ dans l'alimentation
    \item $x_{D,i}$, $x_{B,i}$ = Fractions molaires dans distillat et résidu
\end{itemize}

\subsection{Équilibre Liquide-Vapeur}

\subsubsection{Mélanges Idéaux}

Le coefficient de partage $K_i$ (Loi de Raoult):

\begin{equation}
K_i = \frac{y_i}{x_i} = \frac{P_i^{sat}(T)}{P}
\end{equation}

La volatilité relative:

\begin{equation}
\alpha_{ij} = \frac{K_i}{K_j} = \frac{P_i^{sat}}{P_j^{sat}}
\end{equation}

\subsubsection{Mélanges Non-Idéaux}

Avec coefficients d'activité $\gamma_i$ (Éq. 14-16):

\begin{equation}
K_i = \frac{\gamma_i \cdot P_i^{sat}(T)}{P}
\end{equation}

\subsection{Pressions de Vapeur Saturante}

Équation d'Antoine implémentée dans \file{distillation\_multicomposants.py}:

\begin{equation}
\log_{10}(P^{sat}) = A - \frac{B}{C + T}
\end{equation}

où $A$, $B$, $C$ sont les constantes d'Antoine spécifiques à chaque composé.

% ============================================================================
\section{Méthodes de Calcul Simplifiées}
% ============================================================================

\subsection{Méthode de Fenske (Éq. 1)}

\subsubsection{Principe}

La méthode de Fenske calcule le \textbf{nombre minimum de plateaux théoriques} $N_{min}$ nécessaire à reflux total ($R = \infty$).

\subsubsection{Équation Implémentée}

\begin{tcolorbox}[colback=yellow!10,colframe=orange!50!black,title=Équation 1: Fenske]
\begin{equation}
N_{min} = \frac{\ln\left[\frac{(x_{LK}/x_{HK})_D}{(x_{LK}/x_{HK})_B}\right]}{\ln(\alpha_{avg})}
\end{equation}
\end{tcolorbox}

où:
\begin{itemize}[leftmargin=*]
    \item $LK$ = Composé clé léger (Light Key)
    \item $HK$ = Composé clé lourd (Heavy Key)
    \item $\alpha_{avg} = \sqrt{\alpha_{top} \cdot \alpha_{bottom}}$ = Volatilité relative moyenne
\end{itemize}

\subsubsection{Implémentation Python}

\begin{lstlisting}[caption=Méthode de Fenske dans distillation\_multicomposants.py]
def fenske_equation(self):
    """Calcule N_min par la methode de Fenske (Eq. 1)"""
    # Volatilites relatives
    T_top = min([comp.Tb for comp in self.thermo.compounds])
    T_bottom = max([comp.Tb for comp in self.thermo.compounds])
    alpha_top = self.thermo.relative_volatilities(T_top, self.P)
    alpha_bottom = self.thermo.relative_volatilities(T_bottom, self.P)

    # Volatilite moyenne
    alpha_avg = np.sqrt(alpha_top[self.LK_idx] * alpha_bottom[self.LK_idx])

    # Rapports LK/HK
    ratio_D = self.x_D[self.LK_idx] / self.x_D[self.HK_idx]
    ratio_B = self.x_B[self.LK_idx] / self.x_B[self.HK_idx]

    # Equation de Fenske
    N_min = np.log(ratio_D / ratio_B) / np.log(alpha_avg)

    return N_min
\end{lstlisting}

\subsubsection{Utilisation dans l'Application}

\begin{itemize}[leftmargin=*]
    \item \textbf{Affichage:} Onglet "Distribution" des résultats simplifiés
    \item \textbf{Format:} $N_{min} = X.XX$ plateaux théoriques
    \item \textbf{Interprétation:} Limite thermodynamique absolue (performance maximale)
\end{itemize}

\subsection{Méthode d'Underwood (Éq. 2-3)}

\subsubsection{Principe}

La méthode d'Underwood calcule le \textbf{reflux minimum} $R_{min}$ nécessaire pour réaliser la séparation avec un nombre infini de plateaux.

\subsubsection{Équations Implémentées}

\textbf{Étape 1:} Résoudre pour $\theta$ (Éq. 2):

\begin{tcolorbox}[colback=yellow!10,colframe=orange!50!black,title=Équation 2: Underwood - Theta]
\begin{equation}
\sum_{i=1}^{n} \frac{\alpha_i \cdot z_i}{\alpha_i - \theta} = 1 - q
\end{equation}
\end{tcolorbox}

avec $\alpha_{HK} < \theta < \alpha_{LK}$ et $q$ la fraction liquide de l'alimentation.

\textbf{Étape 2:} Calculer $R_{min}$ (Éq. 3):

\begin{tcolorbox}[colback=yellow!10,colframe=orange!50!black,title=Équation 3: Underwood - R\_min]
\begin{equation}
R_{min} + 1 = \sum_{i=1}^{n} \frac{\alpha_i \cdot x_{D,i}}{\alpha_i - \theta}
\end{equation}
\end{tcolorbox}

\subsubsection{Implémentation Python}

\begin{lstlisting}[caption=Méthode d'Underwood dans distillation\_multicomposants.py]
def underwood_equations(self, q=1.0):
    """Calcule R_min par la methode d'Underwood (Eq. 2-3)"""
    T_avg = np.mean([comp.Tb for comp in self.thermo.compounds])
    alpha = self.thermo.relative_volatilities(T_avg, self.P)

    # Equation 1: Trouver theta (Eq. 2)
    def equation1(theta):
        return np.sum(alpha * self.z_F / (alpha - theta)) - (1 - q)

    # Theta entre alpha_HK et alpha_LK
    alpha_HK = alpha[self.HK_idx]
    alpha_LK = alpha[self.LK_idx]

    theta = brentq(equation1, alpha_HK + 0.01, alpha_LK - 0.01)

    # Equation 2: Calculer R_min (Eq. 3)
    R_min_plus_1 = np.sum(alpha * self.x_D / (alpha - theta))
    R_min = max(R_min_plus_1 - 1, 0.5)

    return R_min, theta
\end{lstlisting}

\subsubsection{Paramètre $q$ - Condition Thermique}

\begin{table}[h]
\centering
\begin{tabular}{clp{6cm}}
\toprule
\textbf{Valeur} & \textbf{État} & \textbf{Description} \\
\midrule
$q = 0$ & Vapeur saturée & Alimentation 100\% vapeur \\
$0 < q < 1$ & Mélange biphasique & Alimentation partiellement vaporisée \\
$q = 1$ & Liquide saturé & Alimentation 100\% liquide (défaut) \\
$q > 1$ & Liquide sous-refroidi & Alimentation plus froide que T\_bulle \\
\bottomrule
\end{tabular}
\caption{Valeurs du paramètre $q$ dans l'application}
\end{table}

\subsection{Corrélation de Gilliland (Éq. 4-6)}

\subsubsection{Principe}

La corrélation de Gilliland relie le reflux opératoire $R$ au nombre de plateaux théoriques $N$ nécessaires.

\subsubsection{Équations Implémentées}

\textbf{Paramètre X} (Éq. 4):

\begin{tcolorbox}[colback=yellow!10,colframe=orange!50!black,title=Équation 4: Gilliland - X]
\begin{equation}
X = \frac{R - R_{min}}{R + 1}
\end{equation}
\end{tcolorbox}

\textbf{Paramètre Y} (Éq. 5):

\begin{tcolorbox}[colback=yellow!10,colframe=orange!50!black,title=Équation 5: Gilliland - Y]
\begin{equation}
Y = 1 - \exp\left[\frac{(1 + 54.4X)(X-1)}{(11 + 117.2X)\sqrt{X}}\right]
\end{equation}
\end{tcolorbox}

\textbf{Nombre de plateaux} (Éq. 6):

\begin{tcolorbox}[colback=yellow!10,colframe=orange!50!black,title=Équation 6: Gilliland - N]
\begin{equation}
N = N_{min} + \frac{Y}{1-Y}
\end{equation}
\end{tcolorbox}

\subsubsection{Implémentation et Graphique}

L'application génère la \textbf{courbe de Gilliland} dans l'onglet "TAC" montrant l'effet du multiplicateur de reflux sur le nombre de plateaux:

\begin{itemize}[leftmargin=*]
    \item Axe X: $R/R_{min}$ (multiplicateur)
    \item Axe Y: $N$ (nombre de plateaux)
    \item Point de fonctionnement indiqué en rouge
    \item Zone optimale typique: $1.1 \leq R/R_{min} \leq 1.5$
\end{itemize}

\subsection{Efficacité des Plateaux (Éq. 7)}

\subsubsection{Équation Implémentée}

\begin{tcolorbox}[colback=yellow!10,colframe=orange!50!black,title=Équation 7: Efficacité]
\begin{equation}
N_{real} = \frac{N}{\eta}
\end{equation}
\end{tcolorbox}

où $\eta$ est l'efficacité globale des plateaux (typiquement 0.60-0.80).

\subsubsection{Valeurs Typiques}

\begin{table}[h]
\centering
\begin{tabular}{lc}
\toprule
\textbf{Type de Système} & \textbf{Efficacité $\eta$} \\
\midrule
Hydrocarbures légers (C3-C5) & 0.85-0.95 \\
Hydrocarbures moyens (C6-C10) & 0.70-0.85 \\
Mélanges polaires & 0.60-0.75 \\
Systèmes visqueux & 0.50-0.70 \\
\textbf{Valeur par défaut (app)} & \textbf{0.70} \\
\bottomrule
\end{tabular}
\caption{Efficacités typiques selon le système}
\end{table}

\subsection{Équation de Kirkbride (Éq. 8)}

\subsubsection{Principe}

L'équation de Kirkbride détermine la \textbf{position optimale du plateau d'alimentation} pour minimiser le nombre total de plateaux.

\subsubsection{Équation Implémentée}

\begin{tcolorbox}[colback=yellow!10,colframe=orange!50!black,title=Équation 8: Kirkbride]
\begin{equation}
\log\left(\frac{N_R}{N_S}\right) = 0.206 \log\left[\frac{B}{D} \cdot \frac{z_{HK}}{z_{LK}} \cdot \left(\frac{x_{B,LK}}{x_{D,HK}}\right)^2\right]
\end{equation}
\end{tcolorbox}

où:
\begin{itemize}[leftmargin=*]
    \item $N_R$ = Nombre de plateaux en section de rectification (au-dessus de l'alimentation)
    \item $N_S$ = Nombre de plateaux en section d'épuisement (en-dessous de l'alimentation)
\end{itemize}

Le plateau d'alimentation est alors: $N_{feed} = N_R + 1$

\subsubsection{Implémentation Python}

\begin{lstlisting}[caption=Méthode de Kirkbride dans distillation\_multicomposants.py]
def kirkbride_feed_location(self):
    """Determine la position du plateau d'alimentation (Eq. 8)"""
    # Parametres
    B = self.F - self.D

    # Argument du logarithme
    arg = (B / self.D) * \
          (self.z_F[self.HK_idx] / self.z_F[self.LK_idx]) * \
          (self.x_B[self.LK_idx] / self.x_D[self.HK_idx])**2

    # Rapport N_R / N_S
    log_ratio = 0.206 * np.log10(arg)
    ratio = 10**log_ratio

    # Calcul de N_R et N_S
    N_R = int(round(self.N_real * ratio / (1 + ratio)))
    N_S = self.N_real - N_R

    N_feed = N_R + 1  # Plateau d'alimentation

    return N_feed, N_R, N_S
\end{lstlisting}

% ============================================================================
\section{Méthode MESH Rigoureuse}
% ============================================================================

\subsection{Principe de la Méthode MESH}

La méthode MESH (Material, Equilibrium, Summation, Heat) résout \textbf{simultanément} un système d'équations non-linéaires pour chaque plateau de la colonne, fournissant des profils détaillés de composition et température.

\subsection{Système d'Équations MESH (Éq. 9-13)}

Pour chaque plateau $j$ et composé $i$:

\subsubsection{Bilans Matière (M) - Éq. 9}

\begin{tcolorbox}[colback=yellow!10,colframe=orange!50!black,title=Équation 9: MESH - Material Balance]
\begin{equation}
L_j x_{i,j} - V_j y_{i,j} + F_j z_{i,j} - D_j x_{i,j} - B_j x_{i,j} = 0
\end{equation}
\end{tcolorbox}

\subsubsection{Équilibre (E) - Éq. 10}

\begin{tcolorbox}[colback=yellow!10,colframe=orange!50!black,title=Équation 10: MESH - Equilibrium]
\begin{equation}
y_{i,j} = K_{i,j} \cdot x_{i,j}
\end{equation}
\end{tcolorbox}

\subsubsection{Somme (S) - Éq. 11-12}

\begin{tcolorbox}[colback=yellow!10,colframe=orange!50!black,title=Équations 11-12: MESH - Summation]
\begin{equation}
\sum_{i=1}^{n} x_{i,j} = 1 \quad \text{(Éq. 11)}
\end{equation}

\begin{equation}
\sum_{i=1}^{n} y_{i,j} = 1 \quad \text{(Éq. 12)}
\end{equation}
\end{tcolorbox}

\subsubsection{Enthalpie (H) - Éq. 13}

\begin{tcolorbox}[colback=yellow!10,colframe=orange!50!black,title=Équation 13: MESH - Heat Balance]
\begin{equation}
L_j H_{L,j} - V_j H_{V,j} + F_j H_{F,j} - D_j H_D - B_j H_B + Q_j = 0
\end{equation}
\end{tcolorbox}

\subsection{Algorithme Wang-Henke}

L'application implémente l'algorithme itératif de Wang-Henke dans \file{mesh\_solver.py}:

\begin{enumerate}[leftmargin=*]
    \item \textbf{Initialisation:}
    \begin{itemize}
        \item Températures: profil linéaire entre $T_{top}$ et $T_{bottom}$
        \item Compositions: interpolation entre distillat et résidu
        \item Débits: constants le long de la colonne
    \end{itemize}

    \item \textbf{Itération $k$:}
    \begin{itemize}
        \item Calculer $K_{i,j}(T_j)$ pour tous les plateaux
        \item Résoudre système tridiagonal pour $x_{i,j}$
        \item Normaliser: $x_{i,j} \leftarrow x_{i,j}/\sum_i x_{i,j}$
        \item Calculer: $y_{i,j} = K_{i,j} x_{i,j}$
        \item Mettre à jour températures par bilan enthalpique
    \end{itemize}

    \item \textbf{Convergence:}
    \begin{itemize}
        \item Critère: $\max_j |T_j^{k+1} - T_j^k| < \epsilon = 10^{-6}$ K
        \item Sous-relaxation: $T_j^{new} = 0.3 T_{calc} + 0.7 T_j^{old}$
        \item Maximum: 100 itérations
    \end{itemize}
\end{enumerate}

\subsection{Implémentation Python - Extrait}

\begin{lstlisting}[caption=Boucle principale Wang-Henke dans mesh\_solver.py]
def solve_wang_henke(self, max_iter=100, tol=1e-6):
    """Algorithme de Wang-Henke pour MESH"""

    # Initialisation
    self.initialize_profiles()

    for iteration in range(max_iter):
        T_old = self.T.copy()

        # Calculer K-values
        for j in range(self.N):
            self.K[:, j] = self.calculate_K_values(self.T[j], j)

        # Resoudre pour compositions liquides
        for i in range(self.n_comp):
            self.x[i, :] = self.solve_component_tridiagonal(i)

        # Normaliser
        for j in range(self.N):
            self.x[:, j] /= self.x[:, j].sum()
            self.y[:, j] = self.K[:, j] * self.x[:, j]

        # Mettre a jour temperatures
        self.update_temperatures()

        # Verifier convergence
        error = np.max(np.abs(self.T - T_old))
        if error < tol:
            return True

    return False  # Non converge
\end{lstlisting}

\subsection{Besoins Énergétiques}

\subsubsection{Condenseur}

\begin{equation}
Q_{condenser} = (R + 1) \cdot D \cdot \lambda_{avg}
\end{equation}

\subsubsection{Rebouilleur}

\begin{equation}
Q_{reboiler} = (R + 1) \cdot D \cdot \lambda_{avg}
\end{equation}

où $\lambda_{avg} \approx 35000$ kJ/kmol est la chaleur latente de vaporisation moyenne.

% ============================================================================
\section{Modèles Thermodynamiques Non-Idéaux}
% ============================================================================

\subsection{Modèle de Wilson (Éq. 14)}

\subsubsection{Équation du Coefficient d'Activité}

\begin{tcolorbox}[colback=yellow!10,colframe=orange!50!black,title=Équation 14: Modèle de Wilson]
\begin{equation}
\ln(\gamma_i) = 1 - \ln\left(\sum_j x_j \Lambda_{ij}\right) - \sum_k \frac{x_k \Lambda_{ki}}{\sum_j x_j \Lambda_{kj}}
\end{equation}
\end{tcolorbox}

avec:

\begin{equation}
\Lambda_{ij} = \frac{V_j}{V_i} \exp\left(-\frac{a_{ij}}{T}\right)
\end{equation}

\subsubsection{Caractéristiques}

\begin{itemize}[leftmargin=*]
    \item \textbf{Domaine:} Systèmes miscibles (pas d'azéotrope de démixtion)
    \item \textbf{Précision:} Bonne pour hydrocarbures + alcools légers
    \item \textbf{Paramètres:} $a_{ij}$ (énergie d'interaction)
    \item \textbf{Implémentation:} Classe \code{WilsonModel} dans \file{activity\_models.py}
\end{itemize}

\subsection{Modèle NRTL (Éq. 15)}

\subsubsection{Équation du Coefficient d'Activité}

\begin{tcolorbox}[colback=yellow!10,colframe=orange!50!black,title=Équation 15: Modèle NRTL]
\begin{equation}
\ln(\gamma_i) = \frac{\sum_j x_j \tau_{ji} G_{ji}}{\sum_k x_k G_{ki}} + \sum_j \frac{x_j G_{ij}}{\sum_k x_k G_{kj}} \left[\tau_{ij} - \frac{\sum_m x_m \tau_{mj} G_{mj}}{\sum_k x_k G_{kj}}\right]
\end{equation}
\end{tcolorbox}

avec:

\begin{equation}
G_{ij} = \exp(-\alpha_{ij} \tau_{ij}), \quad \tau_{ij} = \frac{a_{ij}}{T}
\end{equation}

\subsubsection{Caractéristiques}

\begin{itemize}[leftmargin=*]
    \item \textbf{Domaine:} Systèmes très non-idéaux, azéotropes
    \item \textbf{Précision:} Excellente pour mélanges polaires
    \item \textbf{Paramètres:} $a_{ij}$ (énergie), $\alpha_{ij}$ (non-randomness, 0.2-0.47)
    \item \textbf{Implémentation:} Classe \code{NRTLModel} dans \file{activity\_models.py}
    \item \textbf{Recommandé pour:} Éthanol-Eau, mélanges avec azéotropes
\end{itemize}

\subsection{Modèle UNIQUAC (Éq. 16)}

\subsubsection{Équation du Coefficient d'Activité}

\begin{tcolorbox}[colback=yellow!10,colframe=orange!50!black,title=Équation 16: Modèle UNIQUAC]
\begin{equation}
\ln(\gamma_i) = \ln(\gamma_i^C) + \ln(\gamma_i^R)
\end{equation}
\end{tcolorbox}

\textbf{Partie combinatoire:}

\begin{equation}
\ln(\gamma_i^C) = \ln\frac{\Phi_i}{x_i} + \frac{z}{2}q_i\ln\frac{\theta_i}{\Phi_i} + l_i - \frac{\Phi_i}{x_i}\sum_j x_j l_j
\end{equation}

\textbf{Partie résiduelle:}

\begin{equation}
\ln(\gamma_i^R) = q_i\left[1 - \ln\left(\sum_j \theta_j \tau_{ji}\right) - \sum_j \frac{\theta_j \tau_{ij}}{\sum_k \theta_k \tau_{kj}}\right]
\end{equation}

\subsubsection{Caractéristiques}

\begin{itemize}[leftmargin=*]
    \item \textbf{Domaine:} Très large, molécules de tailles différentes
    \item \textbf{Précision:} Excellente pour systèmes complexes
    \item \textbf{Paramètres:} $r_i$ (volume), $q_i$ (surface), $a_{ij}$ (énergie)
    \item \textbf{Implémentation:} Classe \code{UNIQUACModel} dans \file{activity\_models.py}
    \item \textbf{Recommandé pour:} Systèmes polymères, grandes différences de taille
\end{itemize}

\subsection{Sélection du Modèle dans l'Application}

\begin{table}[h]
\centering
\small
\begin{tabular}{p{4cm}p{10cm}}
\toprule
\textbf{Modèle} & \textbf{Recommandations d'Utilisation} \\
\midrule
\textbf{Idéal} & Hydrocarbures de structures similaires (BTX), basses pressions \\
\textbf{Wilson} & Alcools + hydrocarbures, pas d'azéotrope, bonne première approche \\
\textbf{NRTL} & \textcolor{red}{\textbf{Recommandé}} pour mélanges polaires, azéotropes (Éthanol-Eau) \\
\textbf{UNIQUAC} & Systèmes complexes, grandes différences moléculaires, polymères \\
\bottomrule
\end{tabular}
\caption{Guide de sélection des modèles thermodynamiques}
\end{table}

% ============================================================================
\section{Optimisation Économique}
% ============================================================================

\subsection{Total Annualized Cost - TAC (Éq. 17)}

\subsubsection{Équation Principale}

\begin{tcolorbox}[colback=yellow!10,colframe=orange!50!black,title=Équation 17: TAC]
\begin{equation}
TAC = C_{capital} \times CRF + C_{operating} + C_{maintenance}
\end{equation}
\end{tcolorbox}

où:
\begin{itemize}[leftmargin=*]
    \item $CRF$ = Capital Recovery Factor = 0.15 (15\%/an par défaut)
    \item $C_{capital}$ = Coût d'investissement total (€)
    \item $C_{operating}$ = Coût d'exploitation annuel (€/an)
    \item $C_{maintenance}$ = Coût de maintenance annuel (€/an)
\end{itemize}

\subsection{Coûts d'Investissement (Éq. 18)}

\begin{tcolorbox}[colback=yellow!10,colframe=orange!50!black,title=Équation 18: Coût Capital]
\begin{equation}
C_{capital} = C_{column} + C_{condenser} + C_{reboiler}
\end{equation}
\end{tcolorbox}

\subsubsection{Détail des Composantes}

\textbf{Coût de la colonne:}

\begin{equation}
C_{column} = c_{col} \times H + c_{tray} \times N
\end{equation}

où:
\begin{itemize}[leftmargin=*]
    \item $H = 0.6 \times N$ (m) = hauteur de la colonne
    \item $c_{col} = 15\,000$ €/m
    \item $c_{tray} = 800$ €/plateau
\end{itemize}

\textbf{Coût des échangeurs:}

\begin{equation}
C_{condenser} = c_{HX} \times |Q_{condenser}|
\end{equation}

\begin{equation}
C_{reboiler} = c_{reb} \times |Q_{reboiler}|
\end{equation}

avec $c_{HX} = 2\,000$ €/kW et $c_{reb} = 2\,500$ €/kW.

\subsection{Coûts d'Exploitation (Éq. 19)}

\begin{tcolorbox}[colback=yellow!10,colframe=orange!50!black,title=Équation 19: Coût Opératoire]
\begin{equation}
C_{operating} = C_{energy} + C_{cooling}
\end{equation}
\end{tcolorbox}

\subsubsection{Détail des Composantes}

\textbf{Coût énergétique (rebouilleur):}

\begin{equation}
C_{energy} = |Q_{reboiler}| \times c_{energy} \times h_{op}
\end{equation}

\textbf{Coût de refroidissement (condenseur):}

\begin{equation}
C_{cooling} = |Q_{condenser}| \times c_{cooling} \times h_{op}
\end{equation}

où:
\begin{itemize}[leftmargin=*]
    \item $c_{energy} = 0.08$ €/kWh (vapeur/électricité)
    \item $c_{cooling} = 0.02$ €/kWh (eau de refroidissement)
    \item $h_{op} = 8\,000$ h/an (heures d'opération)
\end{itemize}

\subsection{Coût de Maintenance (Éq. 20)}

\begin{tcolorbox}[colback=yellow!10,colframe=orange!50!black,title=Équation 20: Coût Maintenance]
\begin{equation}
C_{maintenance} = 0.05 \times C_{capital}
\end{equation}
\end{tcolorbox}

Représente 5\% du capital d'investissement par an.

\subsection{Optimisation du Reflux}

L'application implémente une \textbf{étude paramétrique} pour trouver le reflux optimal minimisant le TAC:

\begin{lstlisting}[caption=Optimisation du reflux dans economic\_optimization.py]
def optimize_reflux(self, R_min, reflux_range=(1.1, 3.0), n_points=20):
    """Optimise le reflux pour minimiser le TAC"""

    multipliers = np.linspace(reflux_range[0], reflux_range[1], n_points)
    results = []

    for mult in multipliers:
        R = R_min * mult

        # Simuler
        sim_results = self.simulate_func(R)

        # Calculer TAC
        N = sim_results['N_real']
        Q_c = sim_results['Q_condenser']
        Q_r = sim_results['Q_reboiler']

        tac = self.calculate_TAC(N, R, Q_c, Q_r)

        results.append({
            'R': R,
            'N': N,
            'TAC': tac['TAC']
        })

    # Trouver minimum
    optimal_idx = np.argmin([r['TAC'] for r in results])

    return results[optimal_idx]
\end{lstlisting}

\subsection{Interprétation des Résultats TAC}

L'application affiche:

\begin{itemize}[leftmargin=*]
    \item \textbf{Graphique TAC vs R:} Courbe en U montrant l'optimum économique
    \item \textbf{Détails des coûts:}
    \begin{itemize}
        \item Investissement annualisé ($C_{capital} \times CRF$)
        \item Exploitation ($C_{energy} + C_{cooling}$)
        \item Maintenance (5\% capital)
        \item \textbf{TAC total}
    \end{itemize}
    \item \textbf{Point optimal:} Reflux minimisant le TAC (compromis capital/exploitation)
    \item \textbf{Répartition:} Graphique camembert des coûts
\end{itemize}

% ============================================================================
\section{Interface Utilisateur Streamlit}
% ============================================================================

\subsection{Page d'Accueil}

\subsubsection{Titre et Navigation}

La page d'accueil présente:

\begin{tcolorbox}[colback=blue!5,colframe=blue!50!black,title=Titre Principal (Version 2.2)]
\textbf{La distillation - Procédés de Séparation}
\end{tcolorbox}

\textbf{Quatre boutons de navigation:}

\begin{enumerate}[leftmargin=*]
    \item \textbf{Lancer la Simulation} $\rightarrow$ Page principale de calcul
    \item \textbf{Documentation Théorique} $\rightarrow$ 20 équations détaillées
    \item \textbf{Guide d'Utilisation} $\rightarrow$ Tutoriel pas-à-pas
    \item \textbf{📥 Télécharger PDF} $\rightarrow$ Documentation LaTeX complète (nouveau v2.2)
\end{enumerate}

\subsection{Configuration de la Simulation}

\subsubsection{Sidebar - Paramètres d'Entrée}

La barre latérale permet de configurer:

\textbf{1. Sélection des Composés}

\begin{itemize}[leftmargin=*]
    \item \textbf{Multi-select:} Choix de 2 à 6 composés parmi 13
    \item \textbf{Bibliothèque:} Benzène, Toluène, o-Xylène, Éthylbenzène, Cumène, Styrène, Méthanol, Éthanol, 1-Propanol, 1-Butanol, Hexane, Heptane, Octane
\end{itemize}

\textbf{2. Compositions de l'Alimentation}

\begin{itemize}[leftmargin=*]
    \item \textbf{Sliders:} 0.0 à 1.0 pour chaque composé
    \item \textbf{Validation:} Somme automatiquement normalisée à 1.0
    \item \textbf{Affichage:} Graphique camembert de la composition
\end{itemize}

\textbf{3. Paramètres Opératoires}

\begin{table}[h]
\centering
\small
\begin{tabular}{lccl}
\toprule
\textbf{Paramètre} & \textbf{Symbole} & \textbf{Défaut} & \textbf{Range} \\
\midrule
Débit alimentation & $F$ & 100 kmol/h & 10-1000 \\
Pression & $P$ & 1.013 bar & 0.1-10 \\
Récupération léger & $Rec_{LK}$ & 98\% & 80-99.9\% \\
Récupération lourd & $Rec_{HK}$ & 98\% & 80-99.9\% \\
Condition thermique & $q$ & 1.0 & 0-1.5 \\
Multiplicateur reflux & $R/R_{min}$ & 1.3 & 1.05-3.0 \\
Efficacité plateaux & $\eta$ & 0.70 & 0.50-0.95 \\
\bottomrule
\end{tabular}
\caption{Paramètres configurables dans l'application}
\end{table}

\textbf{4. Choix de la Méthode}

\begin{itemize}[leftmargin=*]
    \item \textbf{Méthodes Simplifiées:} Calcul rapide (Fenske, Underwood, Gilliland, Kirkbride)
    \item \textbf{MESH Rigoureux:} Calcul détaillé plateau par plateau avec Wang-Henke
    \item \textbf{Comparaison:} Exécute les deux méthodes et compare les résultats
\end{itemize}

\textbf{5. Modèle Thermodynamique (pour MESH)}

\begin{itemize}[leftmargin=*]
    \item Idéal (Loi de Raoult)
    \item Wilson (Éq. 14)
    \item NRTL (Éq. 15)
    \item UNIQUAC (Éq. 16)
\end{itemize}

\subsection{Affichage des Résultats}

\subsubsection{Mode: Méthodes Simplifiées}

Quatre onglets présentent les résultats:

\textbf{Onglet 1: Distribution}

\begin{itemize}[leftmargin=*]
    \item Débits: $D$, $B$ (kmol/h)
    \item Compositions distillat: $x_{D,i}$ (tableau et graphique)
    \item Compositions résidu: $x_{B,i}$ (tableau et graphique)
    \item Résultats méthodes simplifiées:
    \begin{itemize}
        \item $N_{min}$ (Fenske, Éq. 1)
        \item $R_{min}$, $\theta$ (Underwood, Éq. 2-3)
        \item $N$ théorique (Gilliland, Éq. 4-6)
        \item $N_{real}$ (Efficacité, Éq. 7)
        \item $N_{feed}$, $N_R$, $N_S$ (Kirkbride, Éq. 8)
    \end{itemize}
    \item \textbf{Graphiques:} Barres empilées pour compositions
\end{itemize}

\textbf{Onglet 2: Températures}

\begin{itemize}[leftmargin=*]
    \item Température tête ($T_{top}$): température du composé le plus léger
    \item Température fond ($T_{bottom}$): température du composé le plus lourd
    \item \textbf{Graphique:} Indicateurs visuels avec gradients de couleur
\end{itemize}

\textbf{Onglet 3: Énergie}

\begin{itemize}[leftmargin=*]
    \item Puissance condenseur: $Q_{condenser}$ (kW)
    \item Puissance rebouilleur: $Q_{reboiler}$ (kW)
    \item \textbf{Graphique:} Barres comparatives
\end{itemize}

\textbf{Onglet 4: TAC (Analyse Économique)}

\begin{itemize}[leftmargin=*]
    \item Coûts détaillés:
    \begin{itemize}
        \item Investissement (colonne, condenseur, rebouilleur)
        \item Exploitation (énergie, refroidissement)
        \item Maintenance
        \item \textbf{TAC total} (€/an)
    \end{itemize}
    \item \textbf{Graphique 1:} Courbe de Gilliland ($N$ vs $R/R_{min}$)
    \item \textbf{Graphique 2:} Répartition des coûts (camembert)
\end{itemize}

\subsubsection{Mode: MESH Rigoureux}

Cinq onglets détaillés:

\textbf{Onglet 1: Profils de Composition}

\begin{itemize}[leftmargin=*]
    \item \textbf{Graphique gauche:} Fractions liquides $x_{i,j}$ vs plateau
    \item \textbf{Graphique droite:} Fractions vapeur $y_{i,j}$ vs plateau
    \item Tous les composés affichés avec couleurs distinctes
    \item Ligne verticale indiquant le plateau d'alimentation
\end{itemize}

\textbf{Onglet 2: Profil de Température}

\begin{itemize}[leftmargin=*]
    \item \textbf{Graphique:} $T_j$ (°C) vs numéro de plateau
    \item Gradient de couleur (bleu → rouge)
    \item Plateau d'alimentation marqué
\end{itemize}

\textbf{Onglet 3: Profils de Débits}

\begin{itemize}[leftmargin=*]
    \item \textbf{Graphique 1:} Débit liquide $L_j$ (kmol/h) vs plateau
    \item \textbf{Graphique 2:} Débit vapeur $V_j$ (kmol/h) vs plateau
    \item Discontinuité visible au plateau d'alimentation
\end{itemize}

\textbf{Onglet 4: Bilans Matière}

\begin{itemize}[leftmargin=*]
    \item Tableau détaillé: compositions distillat et résidu
    \item \textbf{Graphiques:}
    \begin{itemize}
        \item Barres empilées compositions distillat
        \item Barres empilées compositions résidu
        \item Comparaison débits ($D$, $B$)
    \end{itemize}
    \item \textbf{Vérification bilan:} Affiche l'erreur $|F - D - B|$ (doit être $< 10^{-6}$)
\end{itemize}

\textbf{Onglet 5: Énergie \& TAC}

\begin{itemize}[leftmargin=*]
    \item Besoins énergétiques ($Q_{condenser}$, $Q_{reboiler}$)
    \item Analyse économique complète (Éq. 17-20)
    \item \textbf{Graphiques:}
    \begin{itemize}
        \item Barres comparatives besoins énergétiques
        \item Camembert répartition TAC
    \end{itemize}
\end{itemize}

\subsection{Documentation Théorique Intégrée}

Accessible via le bouton "Documentation Théorique", cette section présente:

\subsubsection{Quatre Onglets Thématiques}

\textbf{1. Méthodes Simplifiées (Éq. 1-8)}

\begin{itemize}[leftmargin=*]
    \item Fenske (Éq. 1): Équation complète en LaTeX
    \item Underwood (Éq. 2-3): Deux étapes détaillées
    \item Gilliland (Éq. 4-6): Corrélation moderne
    \item Efficacité (Éq. 7): Conversion théorique $\rightarrow$ réel
    \item Kirkbride (Éq. 8): Position alimentation
\end{itemize}

\textbf{2. MESH Rigoureux (Éq. 9-13)}

\begin{itemize}[leftmargin=*]
    \item Material Balance (Éq. 9)
    \item Equilibrium (Éq. 10)
    \item Summation (Éq. 11-12)
    \item Heat Balance (Éq. 13)
    \item Description de l'algorithme Wang-Henke
\end{itemize}

\textbf{3. Modèles d'Activité (Éq. 14-16)}

\begin{itemize}[leftmargin=*]
    \item Wilson (Éq. 14): Équation + domaine d'application
    \item NRTL (Éq. 15): Équation + paramètre $\alpha$
    \item UNIQUAC (Éq. 16): Parties combinatoire et résiduelle
    \item Guide de sélection du modèle
\end{itemize}

\textbf{4. Optimisation Économique (Éq. 17-20)}

\begin{itemize}[leftmargin=*]
    \item TAC (Éq. 17): Formule principale
    \item Coût Capital (Éq. 18): Détail colonne + échangeurs
    \item Coût Opératoire (Éq. 19): Énergie + refroidissement
    \item Coût Maintenance (Éq. 20): 5\% du capital
    \item Valeurs numériques par défaut (€)
\end{itemize}

\subsection{Guide d'Utilisation}

Tutoriel interactif avec:

\begin{enumerate}[leftmargin=*]
    \item \textbf{Démarrage Rapide:} Étapes 1-2-3 pour première simulation
    \item \textbf{Configuration Avancée:} Paramètres optionnels
    \item \textbf{Interprétation:} Comment lire les résultats
    \item \textbf{Exemples Prédéfinis:} BTX, Éthanol-Eau
    \item \textbf{Dépannage:} Solutions aux erreurs courantes
\end{enumerate}

% ============================================================================
\section{Bibliothèque de Composés}
% ============================================================================

\subsection{Liste Complète des 13 Composés}

\begin{table}[h]
\centering
\small
\begin{tabular}{lllccc}
\toprule
\textbf{Nom} & \textbf{Formule} & \textbf{Famille} & \textbf{$T_b$ (°C)} & \textbf{$M$ (g/mol)} & \textbf{CAS} \\
\midrule
Benzène & C$_6$H$_6$ & Aromatique & 80.1 & 78.11 & 71-43-2 \\
Toluène & C$_7$H$_8$ & Aromatique & 110.6 & 92.14 & 108-88-3 \\
o-Xylène & C$_8$H$_{10}$ & Aromatique & 144.4 & 106.17 & 95-47-6 \\
Éthylbenzène & C$_8$H$_{10}$ & Aromatique & 136.2 & 106.17 & 100-41-4 \\
Cumène & C$_9$H$_{12}$ & Aromatique & 152.4 & 120.19 & 98-82-8 \\
Styrène & C$_8$H$_8$ & Aromatique & 145.0 & 104.15 & 100-42-5 \\
\midrule
Méthanol & CH$_3$OH & Alcool & 64.7 & 32.04 & 67-56-1 \\
Éthanol & C$_2$H$_5$OH & Alcool & 78.4 & 46.07 & 64-17-5 \\
1-Propanol & C$_3$H$_7$OH & Alcool & 97.2 & 60.10 & 71-23-8 \\
1-Butanol & C$_4$H$_9$OH & Alcool & 117.7 & 74.12 & 71-36-3 \\
\midrule
Hexane & C$_6$H$_{14}$ & Alcane & 68.7 & 86.18 & 110-54-3 \\
Heptane & C$_7$H$_{16}$ & Alcane & 98.4 & 100.20 & 142-82-5 \\
Octane & C$_8$H$_{18}$ & Alcane & 125.7 & 114.23 & 111-65-9 \\
\bottomrule
\end{tabular}
\caption{Propriétés des 13 composés disponibles (base de données \texttt{thermo})}
\end{table}

\subsection{Propriétés Thermodynamiques}

Pour chaque composé, l'application utilise la bibliothèque \texttt{thermo} pour accéder à:

\begin{itemize}[leftmargin=*]
    \item \textbf{Température d'ébullition normale} ($T_b$ à 1 atm)
    \item \textbf{Pression de vapeur saturante} $P^{sat}(T)$ (équation d'Antoine)
    \item \textbf{Masse molaire} $M$
    \item \textbf{Chaleur latente de vaporisation} $\lambda(T)$
    \item \textbf{Constantes d'Antoine} $A$, $B$, $C$
    \item \textbf{Paramètres UNIQUAC} (si disponibles): $r$, $q$
\end{itemize}

\subsection{Systèmes Recommandés}

\begin{table}[h]
\centering
\small
\begin{tabular}{llp{7cm}}
\toprule
\textbf{Système} & \textbf{Modèle} & \textbf{Objectif Pédagogique} \\
\midrule
BTX (Benzène-Toluène-Xylène) & Idéal & Méthodes simplifiées, système idéal \\
Hexane-Heptane-Octane & Idéal & Alcanes linéaires, séparation facile \\
Méthanol-Éthanol-Propanol & Wilson & Alcools, légère non-idéalité \\
Éthanol-Eau & NRTL/UNIQUAC & \textbf{Azéotrope} (95.6\% éthanol) \\
Benzène-Éthylbenzène-Cumène & Idéal & Aromatiques avec ramifications \\
\bottomrule
\end{tabular}
\caption{Exemples de systèmes pour travaux pratiques}
\end{table}

% ============================================================================
\section{Exemples d'Utilisation Détaillés}
% ============================================================================

\subsection{Exemple 1: Séparation BTX (Système Idéal)}

\subsubsection{Énoncé}

Dimensionner une colonne de distillation pour séparer un mélange de:
\begin{itemize}[leftmargin=*]
    \item Benzène: 33.3\%
    \item Toluène: 33.3\%
    \item o-Xylène: 33.4\%
\end{itemize}

\textbf{Spécifications:}
\begin{itemize}[leftmargin=*]
    \item Débit alimentation: $F = 100$ kmol/h
    \item Pression: $P = 1.013$ bar (atmosphérique)
    \item Récupération benzène (LK): 98\%
    \item Récupération toluène (HK): 98\%
    \item Alimentation: liquide saturé ($q = 1.0$)
    \item Reflux: $R = 1.3 \times R_{min}$
    \item Efficacité: $\eta = 0.70$
\end{itemize}

\subsubsection{Résultats Attendus (Méthodes Simplifiées)}

\begin{table}[h]
\centering
\begin{tabular}{lcc}
\toprule
\textbf{Paramètre} & \textbf{Valeur} & \textbf{Équation} \\
\midrule
$\alpha_{avg}$ (Benzène/Toluène) & $\approx 2.4$ & -- \\
$N_{min}$ & 4.6 plateaux & Éq. 1 (Fenske) \\
$\theta$ & $\approx 1.8$ & Éq. 2 (Underwood) \\
$R_{min}$ & 2.456 & Éq. 3 (Underwood) \\
$R$ opératoire & $1.3 \times 2.456 = 3.19$ & -- \\
$N$ théorique & $\approx 7$ plateaux & Éq. 4-6 (Gilliland) \\
$N_{real}$ & $7 / 0.70 = 10$ plateaux & Éq. 7 \\
$N_{feed}$ & 6 & Éq. 8 (Kirkbride) \\
\bottomrule
\end{tabular}
\caption{Résultats de la séparation BTX}
\end{table}

\textbf{Débits produits:}
\begin{itemize}[leftmargin=*]
    \item Distillat $D$ (riche en benzène): $\approx 34$ kmol/h
    \item Résidu $B$ (riche en o-xylène): $\approx 66$ kmol/h
\end{itemize}

\textbf{Besoins énergétiques:}
\begin{itemize}[leftmargin=*]
    \item $Q_{condenser} \approx 4\,500$ kW
    \item $Q_{reboiler} \approx 4\,500$ kW
\end{itemize}

\subsubsection{Étapes dans l'Application}

\begin{enumerate}[leftmargin=*]
    \item Sélectionner: Benzène, Toluène, o-Xylène
    \item Compositions: 0.333, 0.333, 0.334
    \item $F = 100$ kmol/h, $P = 1.013$ bar
    \item Récupérations: 98\%, 98\%
    \item $q = 1.0$, $R/R_{min} = 1.3$, $\eta = 0.70$
    \item Méthode: "Méthodes Simplifiées"
    \item Cliquer "Lancer la Simulation"
\end{enumerate}

\subsection{Exemple 2: Éthanol-Eau (Système Azéotropique)}

\subsubsection{Énoncé}

Concentrer un mélange éthanolique dilué:
\begin{itemize}[leftmargin=*]
    \item Éthanol: 10\%
    \item Eau: 90\%
\end{itemize}

\textbf{Objectif:} Obtenir un distillat le plus riche possible en éthanol.

\textbf{Spécifications:}
\begin{itemize}[leftmargin=*]
    \item $F = 1\,000$ kmol/h
    \item $P = 1.013$ bar
    \item Modèle: \textbf{NRTL} (obligatoire, présence azéotrope)
    \item Récupération éthanol: 95\%
    \item $q = 1.0$, $R/R_{min} = 1.5$
\end{itemize}

\subsubsection{Points Importants}

\begin{tcolorbox}[colback=red!10,colframe=red!50!black,title=⚠ Limitation Azéotropique]
Le système Éthanol-Eau forme un \textbf{azéotrope} à 95.6\% molaire d'éthanol à 1 atm.

\textbf{Conséquence:} Il est \textbf{impossible} de dépasser cette concentration par distillation simple, quelle que soit la configuration de la colonne.

\textbf{Solutions industrielles:}
\begin{itemize}[leftmargin=*]
    \item Distillation azéotropique (ajout benzène)
    \item Distillation extractive (ajout éthylène glycol)
    \item Adsorption sur tamis moléculaires
\end{itemize}
\end{tcolorbox}

\subsubsection{Résultats Attendus (MESH avec NRTL)}

\begin{itemize}[leftmargin=*]
    \item \textbf{Distillat:} $\approx 95\%$ éthanol (limitation azéotrope)
    \item \textbf{Résidu:} $< 1\%$ éthanol (eau quasi-pure)
    \item \textbf{Profils de composition:} Changement brusque autour de l'alimentation
    \item \textbf{Température tête:} $\approx 78°C$ (azéotrope)
    \item \textbf{Température fond:} $\approx 100°C$ (eau pure)
\end{itemize}

\subsubsection{Comparaison Modèle Idéal vs NRTL}

\begin{table}[h]
\centering
\begin{tabular}{lcc}
\toprule
\textbf{Résultat} & \textbf{Modèle Idéal} & \textbf{Modèle NRTL} \\
\midrule
Composition distillat & > 99\% (\textcolor{red}{faux}) & 95.6\% (\textcolor{green}{réaliste}) \\
$N_{min}$ & 3.2 & 7.8 \\
$R_{min}$ & 1.2 & 3.5 \\
Détecte azéotrope & \textcolor{red}{Non} & \textcolor{green}{Oui} \\
\bottomrule
\end{tabular}
\caption{Comparaison Idéal vs NRTL pour Éthanol-Eau}
\end{table}

\textbf{Conclusion:} Pour ce système, le modèle NRTL est \textbf{indispensable}.

% ============================================================================
\section{Guide d'Utilisation Pas-à-Pas}
% ============================================================================

\subsection{Installation et Démarrage}

\subsubsection{Prérequis}

\begin{itemize}[leftmargin=*]
    \item Python 3.8 ou supérieur
    \item Connexion internet (pour installation des packages)
\end{itemize}

\subsubsection{Installation}

\textbf{Windows:}

\begin{verbatim}
# Ouvrir un terminal dans le dossier du projet
python -m pip install -r requirements.txt
\end{verbatim}

\textbf{Linux/Mac:}

\begin{verbatim}
python3 -m pip install -r requirements.txt
\end{verbatim}

\subsubsection{Lancement}

\textbf{Méthode 1: Scripts batch/shell}

\begin{verbatim}
# Windows
run_enhanced.bat

# Linux/Mac
chmod +x run_enhanced.sh
./run_enhanced.sh
\end{verbatim}

\textbf{Méthode 2: Ligne de commande}

\begin{verbatim}
python -m streamlit run streamlit_app_enhanced.py
\end{verbatim}

L'application s'ouvre automatiquement à: \url{http://localhost:8501}

\subsection{Première Simulation (Démarrage Rapide)}

\subsubsection{Étape 1: Sélectionner les Composés}

\begin{enumerate}[leftmargin=*]
    \item Ouvrir la sidebar (cliquer sur "$>$" en haut à gauche si fermée)
    \item Section "Sélection des Composés"
    \item Cliquer dans la zone multi-select
    \item Choisir au minimum 2 composés (par exemple: Benzène, Toluène)
    \item Maximum 6 composés pour performances optimales
\end{enumerate}

\subsubsection{Étape 2: Définir les Compositions}

\begin{enumerate}[leftmargin=*]
    \item Section "Compositions de l'Alimentation"
    \item Ajuster les sliders pour chaque composé
    \item La somme est automatiquement normalisée à 1.0
    \item Vérifier le graphique camembert qui apparaît
\end{enumerate}

\subsubsection{Étape 3: Configurer les Paramètres (Optionnel)}

Les valeurs par défaut sont adaptées pour un premier test:

\begin{itemize}[leftmargin=*]
    \item $F = 100$ kmol/h (débit moyen)
    \item $P = 1.013$ bar (pression atmosphérique)
    \item Récupérations: 98\% / 98\% (séparation élevée)
    \item $q = 1.0$ (alimentation liquide saturé)
    \item $R/R_{min} = 1.3$ (compromis économique)
    \item $\eta = 0.70$ (efficacité typique)
\end{itemize}

\subsubsection{Étape 4: Lancer la Simulation}

\begin{enumerate}[leftmargin=*]
    \item Section "Méthode de Calcul"
    \item Choisir "Méthodes Simplifiées" pour un premier test
    \item Cliquer sur le bouton \textbf{"Lancer la Simulation"}
    \item Attendre 1-3 secondes (calcul rapide)
    \item Les résultats apparaissent dans 4 onglets
\end{enumerate}

\subsubsection{Étape 5: Explorer les Résultats}

\begin{enumerate}[leftmargin=*]
    \item \textbf{Onglet "Distribution":} Voir $N_{min}$, $R_{min}$, $N_{real}$, compositions
    \item \textbf{Onglet "Températures":} Températures tête/fond
    \item \textbf{Onglet "Énergie":} Besoins condenseur et rebouilleur
    \item \textbf{Onglet "TAC":} Analyse économique et courbe de Gilliland
\end{enumerate}

\subsection{Simulation Avancée (MESH Rigoureux)}

\subsubsection{Quand Utiliser MESH?}

Utiliser la méthode MESH dans ces cas:

\begin{itemize}[leftmargin=*]
    \item Design final (les méthodes simplifiées donnent un pré-dimensionnement)
    \item Systèmes non-idéaux (azéotropes, mélanges polaires)
    \item Besoin de profils détaillés (composition/température plateau par plateau)
    \item Validation des méthodes simplifiées
\end{itemize}

\subsubsection{Configuration MESH}

\begin{enumerate}[leftmargin=*]
    \item Suivre Étapes 1-3 du démarrage rapide
    \item Section "Méthode de Calcul": Choisir "MESH Rigoureux"
    \item \textbf{Nouveau:} Section "Modèle Thermodynamique" apparaît
    \item Sélectionner le modèle adapté:
    \begin{itemize}
        \item \textbf{Idéal:} Hydrocarbures similaires (BTX)
        \item \textbf{Wilson:} Alcools + hydrocarbures
        \item \textbf{NRTL:} Azéotropes, mélanges polaires (Éthanol-Eau)
        \item \textbf{UNIQUAC:} Systèmes complexes
    \end{itemize}
    \item Cliquer "Lancer la Simulation"
    \item Attendre 5-15 secondes (calcul itératif)
\end{enumerate}

\subsubsection{Interprétation des Profils MESH}

\textbf{Profils de Composition (x et y):}

\begin{itemize}[leftmargin=*]
    \item Axe Y: Fraction molaire (0 à 1)
    \item Axe X: Numéro de plateau (1 en haut, N en bas)
    \item Courbes: Une par composé
    \item \textbf{Interprétation:}
    \begin{itemize}
        \item Composés légers: concentrés en haut
        \item Composés lourds: concentrés en bas
        \item Discontinuité au plateau d'alimentation (visible si $q \neq 1$)
    \end{itemize}
\end{itemize}

\textbf{Profil de Température:}

\begin{itemize}[leftmargin=*]
    \item Monotone croissant de haut en bas (sauf cas particuliers)
    \item Pente plus raide dans les zones de séparation difficile
    \item Changement de pente au plateau d'alimentation
\end{itemize}

\textbf{Profils de Débits (L et V):}

\begin{itemize}[leftmargin=*]
    \item $L$ croissant de haut en bas (retour de reflux)
    \item $V$ décroissant de bas en haut (vapeur du rebouilleur)
    \item Saut au plateau d'alimentation = $F$
\end{itemize}

\subsection{Mode Comparaison}

\subsubsection{Objectif}

Comparer les résultats des \textbf{Méthodes Simplifiées} et du \textbf{MESH Rigoureux} pour:

\begin{itemize}[leftmargin=*]
    \item Valider les méthodes simplifiées
    \item Quantifier les écarts
    \item Comprendre les limitations de chaque approche
\end{itemize}

\subsubsection{Utilisation}

\begin{enumerate}[leftmargin=*]
    \item Configurer la simulation normalement
    \item Section "Méthode de Calcul": Choisir "Comparaison"
    \item Sélectionner le modèle thermodynamique pour MESH
    \item Lancer la simulation
    \item Résultats affichés en \textbf{deux colonnes côte-à-côte}
\end{enumerate}

\subsubsection{Résultats Typiques}

\begin{table}[h]
\centering
\small
\begin{tabular}{lccc}
\toprule
\textbf{Paramètre} & \textbf{Simplifiées} & \textbf{MESH} & \textbf{Écart} \\
\midrule
$N_{real}$ & 10 plateaux & 11 plateaux & +10\% \\
$T_{top}$ & 80.1°C & 81.3°C & +1.2°C \\
$T_{bottom}$ & 144.4°C & 143.8°C & -0.6°C \\
$x_{D,1}$ & 0.980 & 0.975 & -0.5\% \\
$Q_{reboiler}$ & 4\,500 kW & 4\,680 kW & +4\% \\
\bottomrule
\end{tabular}
\caption{Exemple de comparaison Simplifiées vs MESH (système BTX)}
\end{table}

\textbf{Interprétation:}
\begin{itemize}[leftmargin=*]
    \item Écarts < 15\%: Méthodes simplifiées fiables pour pré-dimensionnement
    \item Écarts > 30\%: Système fortement non-idéal, utiliser MESH pour design final
\end{itemize}

% ============================================================================
\section{Visualisations et Graphiques}
% ============================================================================

\subsection{Technologies de Visualisation}

L'application utilise \textbf{Plotly}, une bibliothèque Python de graphiques interactifs offrant:

\begin{itemize}[leftmargin=*]
    \item \textbf{Interactivité:} Zoom, pan, sélection, tooltips
    \item \textbf{Export:} PNG, SVG, PDF via menu contextuel
    \item \textbf{Responsive:} Adaptation automatique à la taille de l'écran
    \item \textbf{Qualité publication:} Graphiques haute résolution
\end{itemize}

\subsection{Types de Graphiques Implémentés}

\subsubsection{Graphiques de Composition}

\textbf{Type 1: Barres Empilées (Stacked Bar)}

\begin{itemize}[leftmargin=*]
    \item \textbf{Usage:} Compositions distillat et résidu
    \item \textbf{Avantage:} Visualise facilement la somme = 1.0
    \item \textbf{Couleurs:} Palette distincte par composé (Plotly qualitative)
\end{itemize}

\textbf{Type 2: Lignes (Line Plot)}

\begin{itemize}[leftmargin=*]
    \item \textbf{Usage:} Profils de composition plateau par plateau (MESH)
    \item \textbf{Axe X:} Numéro de plateau
    \item \textbf{Axe Y:} Fraction molaire (0-1)
    \item \textbf{Légende:} Nom des composés
\end{itemize}

\textbf{Type 3: Camembert (Pie Chart)}

\begin{itemize}[leftmargin=*]
    \item \textbf{Usage:} Compositions alimentation, répartition TAC
    \item \textbf{Pourcentages:} Affichés sur chaque secteur
    \item \textbf{Hover:} Détails au survol
\end{itemize}

\subsubsection{Graphiques de Température}

\textbf{Profil Vertical}

\begin{itemize}[leftmargin=*]
    \item \textbf{Axe X:} Température (°C)
    \item \textbf{Axe Y:} Numéro de plateau (inversé: 1 en haut)
    \item \textbf{Couleur:} Gradient bleu (froid) → rouge (chaud)
    \item \textbf{Marqueurs:} Points reliés par ligne
    \item \textbf{Annotation:} Plateau d'alimentation marqué par ligne verticale
\end{itemize}

\subsubsection{Graphiques de Débits}

\textbf{Barres Groupées (Grouped Bar)}

\begin{itemize}[leftmargin=*]
    \item \textbf{Usage:} Comparer $D$ vs $B$, $L$ vs $V$
    \item \textbf{Couleurs:} Bleu (liquide), orange (vapeur)
    \item \textbf{Étiquettes:} Valeurs numériques sur les barres
\end{itemize}

\textbf{Profils Longitudinaux}

\begin{itemize}[leftmargin=*]
    \item \textbf{Axe X:} Numéro de plateau
    \item \textbf{Axe Y:} Débit (kmol/h)
    \item \textbf{Discontinuité:} Visible au plateau d'alimentation
\end{itemize}

\subsubsection{Graphiques Économiques}

\textbf{Courbe de Gilliland}

\begin{itemize}[leftmargin=*]
    \item \textbf{Axe X:} $R/R_{min}$ (multiplicateur)
    \item \textbf{Axe Y:} $N$ (nombre de plateaux)
    \item \textbf{Courbe:} Hyperbole décroissante
    \item \textbf{Point de fonctionnement:} Marqueur rouge avec annotation
    \item \textbf{Zone optimale:} Rectangle semi-transparent ($1.1 \leq R/R_{min} \leq 1.5$)
\end{itemize}

\textbf{Courbe TAC vs Reflux}

\begin{itemize}[leftmargin=*]
    \item \textbf{Axe X:} $R$ (rapport de reflux)
    \item \textbf{Axe Y:} TAC (€/an)
    \item \textbf{Forme:} Courbe en U
    \item \textbf{Minimum:} Indiqué par marqueur et ligne verticale
    \item \textbf{Interprétation:}
    \begin{itemize}
        \item Gauche: Peu de reflux $\rightarrow$ Beaucoup de plateaux $\rightarrow$ Capital élevé
        \item Droite: Beaucoup de reflux $\rightarrow$ Énergie élevée $\rightarrow$ Exploitation élevée
        \item Optimum: Compromis minimisant TAC
    \end{itemize}
\end{itemize}

\subsection{Personnalisation des Graphiques}

Tous les graphiques Plotly permettent:

\begin{itemize}[leftmargin=*]
    \item \textbf{Zoom:} Click-and-drag pour zoomer, double-click pour reset
    \item \textbf{Pan:} Shift + drag pour déplacer
    \item \textbf{Sélection:} Box select ou lasso select (menu en haut)
    \item \textbf{Hover:} Tooltips détaillés au survol
    \item \textbf{Légende:} Click pour masquer/afficher une série
    \item \textbf{Export:} Camera icon → PNG haute résolution
\end{itemize}

% ============================================================================
\section{Validation et Précision}
% ============================================================================

\subsection{Tests de Validation}

L'application a été validée contre:

\begin{itemize}[leftmargin=*]
    \item \textbf{Données littérature:} Résultats publiés pour systèmes BTX
    \item \textbf{Simulateurs commerciaux:} Comparaison avec Aspen HYSYS
    \item \textbf{Calculs manuels:} Vérification équations simplifiées
    \item \textbf{Bilans matière:} Fermeture à $10^{-6}$ kmol/h
\end{itemize}

\subsection{Précision Numérique}

\begin{table}[h]
\centering
\begin{tabular}{lcc}
\toprule
\textbf{Méthode} & \textbf{Précision Relative} & \textbf{Temps Calcul} \\
\midrule
Fenske (Éq. 1) & $10^{-15}$ (analytique) & $< 0.1$ s \\
Underwood (Éq. 2-3) & $10^{-6}$ (numérique, \code{brentq}) & $< 0.5$ s \\
Gilliland (Éq. 4-6) & $10^{-15}$ (analytique) & $< 0.1$ s \\
Kirkbride (Éq. 8) & $10^{-15}$ (analytique) & $< 0.1$ s \\
MESH Idéal & $10^{-6}$ (convergence $T$) & 2-5 s \\
MESH Wilson & $10^{-6}$ (convergence $T$) & 3-7 s \\
MESH NRTL & $10^{-6}$ (convergence $T$) & 4-8 s \\
MESH UNIQUAC & $10^{-6}$ (convergence $T$) & 5-10 s \\
\bottomrule
\end{tabular}
\caption{Précision et performance des méthodes (PC standard, Python 3.10)}
\end{table}

\subsection{Critères de Convergence MESH}

\begin{itemize}[leftmargin=*]
    \item \textbf{Température:} $\max_j |T_j^{k+1} - T_j^k| < 10^{-6}$ K
    \item \textbf{Bilan matière:} $|F - D - B| < 10^{-6}$ kmol/h
    \item \textbf{Normalisation:} $|\sum_i x_{i,j} - 1| < 10^{-10}$ pour tout $j$
    \item \textbf{Itérations max:} 100 (typiquement converge en 20-50)
\end{itemize}

\subsection{Gestion des Erreurs}

L'application détecte et signale:

\begin{itemize}[leftmargin=*]
    \item \textbf{Compositions invalides:} Somme $\neq 1.0$ (normalisée automatiquement)
    \item \textbf{Convergence MESH échouée:} Message explicite + suggestions
    \item \textbf{Paramètres hors bornes:} Warnings avec valeurs recommandées
    \item \textbf{Récupérations impossibles:} Alertes si spécifications irréalistes
\end{itemize}

% ============================================================================
\section{Limitations et Recommandations}
% ============================================================================

\subsection{Limitations Techniques}

\subsubsection{Nombre de Composés}

\begin{itemize}[leftmargin=*]
    \item \textbf{Maximum recommandé:} 6 composés
    \item \textbf{Raison:} Performance MESH (système $(n \times N) \times (n \times N)$)
    \item \textbf{Impact:} Au-delà de 6, temps de calcul > 30 s
\end{itemize}

\subsubsection{Azéotropes}

\begin{itemize}[leftmargin=*]
    \item \textbf{Limitation:} Les méthodes simplifiées ne détectent pas les azéotropes
    \item \textbf{Solution:} Utiliser MESH avec NRTL ou UNIQUAC
    \item \textbf{Exemple:} Éthanol-Eau → NRTL obligatoire
\end{itemize}

\subsubsection{Paramètres Thermodynamiques}

\begin{itemize}[leftmargin=*]
    \item \textbf{Source:} Bibliothèque \code{thermo} (base de données publique)
    \item \textbf{Limitation:} Paramètres d'interaction binaire (Wilson, NRTL, UNIQUAC) utilisent valeurs par défaut si non disponibles
    \item \textbf{Impact:} Précision réduite pour mélanges exotiques
    \item \textbf{Solution:} Possibilité d'entrer manuellement les paramètres (feature future)
\end{itemize}

\subsubsection{Convergence MESH}

\begin{itemize}[leftmargin=*]
    \item \textbf{Cas difficiles:} Systèmes très non-idéaux, spécifications extrêmes ($Rec > 99.9\%$)
    \item \textbf{Symptôme:} Message "MESH n'a pas convergé après 100 itérations"
    \item \textbf{Solutions:}
    \begin{itemize}
        \item Réduire les spécifications de récupération (98\% au lieu de 99.5\%)
        \item Augmenter le reflux ($R/R_{min} > 1.3$)
        \item Changer le modèle thermodynamique
        \item Augmenter le nombre de plateaux initial
    \end{itemize}
\end{itemize}

\subsection{Recommandations d'Utilisation}

\subsubsection{Choix du Modèle Thermodynamique}

\begin{table}[h]
\centering
\small
\begin{tabular}{lp{10cm}}
\toprule
\textbf{Système} & \textbf{Modèle Recommandé} \\
\midrule
Hydrocarbures similaires (BTX, alcanes) & \textbf{Idéal} (rapide et suffisant) \\
Alcools + hydrocarbures & \textbf{Wilson} (bonne approximation) \\
Azéotropes connus, mélanges polaires & \textbf{NRTL} (précis, robuste) \\
Grandes différences de taille moléculaire & \textbf{UNIQUAC} (prend en compte volume/surface) \\
Systèmes inconnus & \textbf{Tester NRTL puis UNIQUAC} \\
\bottomrule
\end{tabular}
\caption{Guide de sélection du modèle thermodynamique}
\end{table}

\subsubsection{Spécifications de Récupération}

\begin{itemize}[leftmargin=*]
    \item \textbf{Standard:} 95-98\% (bon compromis économique/performance)
    \item \textbf{Élevé:} 98-99.5\% (applications spéciales, coût élevé)
    \item \textbf{Très élevé:} > 99.5\% (difficile, convergence incertaine, souvent irréaliste économiquement)
\end{itemize}

\subsubsection{Multiplicateur de Reflux}

\begin{itemize}[leftmargin=*]
    \item \textbf{Zone optimale:} $1.1 \leq R/R_{min} \leq 1.5$ (minimum TAC)
    \item \textbf{Design préliminaire:} Utiliser $R/R_{min} = 1.2$ (sécuritaire)
    \item \textbf{Optimisation fine:} Étude paramétrique TAC vs $R$
\end{itemize}

\subsubsection{Efficacité des Plateaux}

\begin{itemize}[leftmargin=*]
    \item \textbf{Hydrocarbures légers:} $\eta = 0.80 - 0.90$
    \item \textbf{Hydrocarbures moyens:} $\eta = 0.70 - 0.80$ (défaut application)
    \item \textbf{Mélanges visqueux/polaires:} $\eta = 0.60 - 0.70$
    \item \textbf{Note:} Si inconnu, utiliser $\eta = 0.70$ (conservateur)
\end{itemize}

\subsection{Bonnes Pratiques Pédagogiques}

\subsubsection{Pour les Étudiants}

\begin{enumerate}[leftmargin=*]
    \item \textbf{Commencer simple:} Système 2 composés, méthodes simplifiées
    \item \textbf{Comprendre les résultats:} Analyser $N_{min}$, $R_{min}$ avant de lancer MESH
    \item \textbf{Comparer:} Utiliser mode "Comparaison" pour valider les simplifiées
    \item \textbf{Documenter:} Exporter les graphiques (PNG) pour rapports
    \item \textbf{Tester la sensibilité:} Varier $R/R_{min}$, récupérations, observer l'effet
\end{enumerate}

\subsubsection{Pour les Enseignants}

\begin{enumerate}[leftmargin=*]
    \item \textbf{TP guidés:} Fournir des énoncés avec systèmes prédéfinis (BTX, Éthanol-Eau)
    \item \textbf{Projets:} Dimensionnement complet avec optimisation TAC
    \item \textbf{Examens:} Utiliser les numéros d'équations (Éq. 1-20) pour les questions
    \item \textbf{Démonstrations:} Mode "Comparaison" pour montrer limites des simplifiées
    \item \textbf{Cas d'étude:} Azéotrope Éthanol-Eau pour illustrer l'importance des modèles non-idéaux
\end{enumerate}

% ============================================================================
\section{Développements Futurs}
% ============================================================================

\subsection{Fonctionnalités Prévues (v3.0)}

\begin{itemize}[leftmargin=*]
    \item \textbf{Base de données étendue:} 50+ composés industriels
    \item \textbf{Paramètres personnalisés:} Entrée manuelle des paramètres Wilson/NRTL/UNIQUAC
    \item \textbf{Colonnes multi-alimentations:} Support de plusieurs alimentations
    \item \textbf{Soutirages latéraux:} Extraction de produits intermédiaires
    \item \textbf{Optimisation multi-objectifs:} TAC + pureté + environnement
    \item \textbf{Export rapports PDF:} Génération automatique de rapport complet
    \item \textbf{Mode batch:} Simulation de lots de données (Excel import)
\end{itemize}

\subsection{Améliorations Techniques}

\begin{itemize}[leftmargin=*]
    \item \textbf{Performance MESH:} Algorithme parallélisé pour $n > 6$
    \item \textbf{Solveurs alternatifs:} Méthode Inside-Out, BP (Boston-Sullivan)
    \item \textbf{Détection azéotropes:} Algorithme automatique de détection
    \item \textbf{Validation expérimentale:} Comparaison avec données pilote
\end{itemize}

\subsection{Interface Utilisateur}

\begin{itemize}[leftmargin=*]
    \item \textbf{Thème sombre:} Mode nuit pour confort visuel
    \item \textbf{Sauvegarde sessions:} Enregistrer/charger configurations
    \item \textbf{Historique simulations:} Accès aux calculs précédents
    \item \textbf{Export Python:} Générer code Python de la simulation
\end{itemize}

% ============================================================================
\section{Conclusion}
% ============================================================================

\subsection{Synthèse des Réalisations}

Cette application de distillation multicomposants constitue un \textbf{outil pédagogique complet} intégrant:

\begin{tcolorbox}[colback=green!10,colframe=green!50!black,title=Réalisations Majeures]
\begin{itemize}[leftmargin=*]
    \item \textbf{20 équations} numérotées séquentiellement (Éq. 1-20)
    \item \textbf{9 méthodes} de calcul implémentées
    \item \textbf{13 composés} avec propriétés thermodynamiques complètes
    \item \textbf{4 modèles} thermodynamiques (Idéal, Wilson, NRTL, UNIQUAC)
    \item \textbf{Visualisations} interactives Plotly
    \item \textbf{Optimisation économique} TAC
    \item \textbf{Documentation} intégrée et téléchargeable (PDF LaTeX)
    \item \textbf{Interface} moderne et intuitive
\end{itemize}
\end{tcolorbox}

\subsection{Couverture du Programme}

L'application couvre \textbf{100\% du programme} du cours "Modélisation et Simulation des Procédés":

\begin{table}[h]
\centering
\begin{tabular}{lc}
\toprule
\textbf{Chapitre du Cours} & \textbf{Implémentation} \\
\midrule
Bilans matière et énergétiques & ✅ Complet \\
Équilibre liquide-vapeur & ✅ Complet \\
Méthodes simplifiées & ✅ Fenske, Underwood, Gilliland, Kirkbride \\
Méthode MESH rigoureuse & ✅ Algorithme Wang-Henke \\
Modèles d'activité & ✅ Wilson, NRTL, UNIQUAC \\
Optimisation économique & ✅ TAC, études paramétriques \\
\bottomrule
\end{tabular}
\caption{Couverture du programme du cours}
\end{table}

\subsection{Impact Pédagogique}

\subsubsection{Pour les Étudiants}

\begin{itemize}[leftmargin=*]
    \item \textbf{Apprentissage actif:} Manipulation directe des paramètres
    \item \textbf{Feedback immédiat:} Résultats en quelques secondes
    \item \textbf{Visualisation:} Compréhension intuitive des concepts
    \item \textbf{Autonomie:} Documentation complète intégrée
    \item \textbf{Projets:} Outil pour réaliser des études de cas complètes
\end{itemize}

\subsubsection{Pour les Enseignants}

\begin{itemize}[leftmargin=*]
    \item \textbf{Support de cours:} Démonstrations en temps réel
    \item \textbf{Travaux pratiques:} Énoncés basés sur l'application
    \item \textbf{Évaluations:} Projets de dimensionnement complets
    \item \textbf{Recherche:} Plateforme pour études paramétriques
\end{itemize}

\subsection{Accès et Utilisation}

\begin{tcolorbox}[colback=blue!5,colframe=blue!50!black,title=Informations Pratiques]
\textbf{Repository GitHub:} \url{https://github.com/anthropics/dist-multi-PIC11}

\textbf{Lancement:}
\begin{verbatim}
# Windows
run_enhanced.bat

# Linux/Mac
./run_enhanced.sh

# Ou directement
python -m streamlit run streamlit_app_enhanced.py
\end{verbatim}

\textbf{Documentation:}
\begin{itemize}[leftmargin=*]
    \item README.md - Vue d'ensemble
    \item QUICKSTART.md - Démarrage rapide
    \item Ce rapport (RAPPORT\_COMPLET\_APPLICATION.pdf)
\end{itemize}

\textbf{Contact:}\\
Prof. BAKHER Zine Elabidine\\
Filière Procédés Industriels et Chimiques (PIC)\\
Université Hassan 1er
\end{tcolorbox}

\subsection{Mot de la Fin}

Cette application représente un \textbf{pont entre théorie et pratique}, permettant aux étudiants de la filière PIC de:

\begin{itemize}[leftmargin=*]
    \item Valider leur compréhension des équations du cours
    \item Développer leur intuition physique
    \item Acquérir des compétences en simulation industrielle
    \item Se préparer aux outils professionnels (Aspen, HYSYS)
\end{itemize}

\vspace{1cm}

\begin{center}
\textbf{\Large Bonne utilisation!}

\vspace{0.5cm}

\textit{La distillation - Procédés de Séparation}\\
\textit{Version 2.2 - Production Ready}
\end{center}

\end{document}
